\begingroup
\renewcommand{\clearpageforsection}{}
\exercisesheader{}
\endgroup

% 7

\eoce{\qt{Penyakit kronis, Bagian I\label{chronic_illness_intro}} 
Pada 2013, Pew Research Foundation melaporkan bahwa ``45\% orang dewasa AS menyatakan 
bahwa mereka hidup dengan satu atau lebih penyakit kronis''.
\footfullcite{data:pewdiagnosis:2013} Namun, nilai ini didasarkan pada suatu sampel, 
sehingga nilai itu sendiri mungkin bukan estimasi yang sempurna bagi parameter populasi yang menjadi 
perhatian. Studi tersebut melaporkan galat baku sekitar 1.2\%, dan model normal 
dapat digunakan secara masuk akal dalam situasi ini. Buatlah interval kepercayaan 95\% bagi 
proporsi orang dewasa AS yang hidup dengan satu atau lebih penyakit kronis. Selain itu, 
tafsirkan interval kepercayaan tersebut dalam konteks studi.
}{}

% 8

\eoce{\qt{Pengguna Twitter dan berita, Bagian I\label{twitter_users_intro}} 
Sebuah jajak pendapat yang dilakukan pada 2013 menemukan bahwa 52\% pengguna Twitter dewasa di AS 
mendapatkan setidaknya sebagian berita dari Twitter.\footfullcite{data:pewtwitternews:2013}. 
Galat baku untuk estimasi ini adalah 2.4\%, dan distribusi normal 
dapat digunakan untuk memodelkan proporsi sampel. Buatlah interval kepercayaan 
99\% bagi proporsi pengguna Twitter dewasa di AS yang mendapatkan sebagian 
berita dari Twitter, dan tafsirkan interval kepercayaan tersebut dalam konteksnya.
}{}

% 9

\eoce{\qt{Penyakit kronis, Bagian II\label{chronic_illness_tf}} Pada 2013, Pew Research Foundation melaporkan bahwa 
``45\% orang dewasa AS menyatakan bahwa mereka hidup dengan satu atau lebih penyakit 
kronis'', dan galat baku untuk estimasi ini adalah 1.2\%. Tentukan apakah setiap 
pernyataan berikut benar atau salah. Berikan penjelasan untuk membenarkan 
setiap jawaban Anda.
\begin{parts}
\item Kita dapat mengatakan dengan pasti bahwa interval kepercayaan dari 
Latihan~\ref{chronic_illness_intro} mencakup persentase sebenarnya dari orang dewasa AS yang 
mengidap penyakit kronis.
\item Jika kita mengulangi studi ini 1,000 kali dan membuat interval kepercayaan 95\% 
untuk setiap studi, maka sekitar 950 dari interval kepercayaan tersebut 
akan mencakup proporsi sebenarnya dari orang dewasa AS yang mengidap penyakit kronis.
\item Jajak pendapat tersebut menyediakan bukti yang signifikan secara statistik (pada 
tingkat $\alpha = 0.05$) bahwa persentase orang dewasa AS yang mengidap 
penyakit kronis berada di bawah 50\%.
\item Karena galat bakunya 1.2\%, hanya 1.2\% orang dalam studi tersebut 
yang menyatakan ketidakpastian tentang jawaban mereka.
\end{parts}
}{}

% 10

\eoce{\qt{Pengguna Twitter dan berita, Bagian II\label{twitter_users_tf}} Sebuah jajak pendapat yang dilakukan pada 2013 menemukan bahwa 52\% 
pengguna Twitter dewasa di AS mendapatkan setidaknya sebagian berita dari Twitter, dan galat baku 
untuk estimasi ini adalah 2.4\%. Tentukan apakah setiap pernyataan berikut 
benar atau salah. Berikan penjelasan untuk membenarkan setiap jawaban Anda.
\begin{parts}
\item Data tersebut menyediakan bukti yang signifikan secara statistik bahwa lebih dari separuh 
pengguna Twitter dewasa di AS mendapatkan sebagian berita melalui Twitter. Gunakan tingkat signifikansi 
$\alpha = 0.01$.
(Bagian ini menggunakan konsep dari Bagian~\ref{hypothesisTesting} dan akan
diperbaiki dalam edisi mendatang.)
\item Karena galat bakunya 2.4\%, kita dapat menyimpulkan bahwa 97.6\% dari seluruh 
pengguna Twitter dewasa di AS disertakan dalam studi tersebut.
\item Jika ingin mengurangi galat baku estimasi, kita seharusnya mengumpulkan 
lebih sedikit data.
\item Jika kita membuat interval kepercayaan 90\% bagi persentase 
pengguna Twitter dewasa di AS yang mendapatkan sebagian berita melalui Twitter, interval kepercayaan ini 
akan lebih lebar daripada interval kepercayaan 99\% yang bersesuaian.
\end{parts}
}{}

% Boundary-clean reader removes a legacy forced page break.

% 11

\eoce{\qt{Menunggu di IGD, Bagian I\label{er_wait_intro_prop_ok}} Seorang administrator rumah sakit 
yang berharap dapat memperbaiki waktu tunggu memutuskan untuk mengestimasi rata-rata waktu 
tunggu di instalasi gawat darurat rumah sakitnya. Ia mengambil sampel acak sederhana 
yang terdiri atas 64 pasien dan menentukan waktu (dalam menit) sejak mereka 
mendaftar di IGD hingga pertama kali diperiksa oleh seorang dokter. Interval kepercayaan 
95\% berdasarkan sampel ini adalah (128 menit, 147 menit), 
yang didasarkan pada model normal untuk rata-rata. Tentukan apakah 
pernyataan-pernyataan berikut benar atau salah, dan jelaskan alasan Anda.
\begin{parts}
\item Kita yakin 95\% bahwa rata-rata waktu tunggu dari 64 pasien instalasi 
gawat darurat ini berada antara 128 dan 147 menit.
\item Kita yakin 95\% bahwa rata-rata waktu tunggu semua pasien di 
instalasi gawat darurat rumah sakit ini berada antara 128 dan 147 menit.
\item 95\% sampel acak memiliki rata-rata sampel antara 128 dan 147 menit.
\item Interval kepercayaan 99\% akan lebih sempit daripada interval kepercayaan 
95\% karena kita perlu lebih yakin terhadap estimasi kita.
\item Batas galatnya adalah 9.5 dan rata-rata sampelnya adalah 137.5.
\item Untuk memperkecil batas galat interval kepercayaan 95\% menjadi 
separuh nilainya saat ini, kita perlu menggandakan ukuran sampel.
(Petunjuk: batas galat bagi rata-rata berskala dengan cara yang sama terhadap ukuran sampel
seperti batas galat bagi proporsi.)
\end{parts}
}{}

% 12

\eoce{\qt{Kesehatan mental\label{mental_health}}
General Social Survey mengajukan pertanyaan:
``Selama berapa hari dalam 30 hari terakhir kondisi 
kesehatan mental Anda, yang mencakup stres, depresi,
dan masalah emosi, tidak baik?"
Berdasarkan tanggapan dari 1,151 penduduk AS,
survei tersebut melaporkan interval kepercayaan 95\% sebesar
3.40 hingga 4.24 hari pada 2010.
\begin{parts}
\item
    Tafsirkan interval ini dalam konteks data.
\item
    Apa arti ``yakin 95\%"? Jelaskan dalam
    konteks penerapannya.
\item
    Misalkan para peneliti berpendapat bahwa tingkat kepercayaan 99\%
    akan lebih sesuai untuk interval ini.
    Apakah interval yang baru ini akan lebih sempit atau lebih lebar daripada
    interval kepercayaan 95\%?
\item
    Jika sebuah survei baru dilakukan terhadap 500 orang Amerika,
    apakah menurut Anda galat baku estimasinya akan
    lebih besar, lebih kecil, atau kurang lebih sama.
\end{parts}
}{}

% 13

\eoce{\qt{Pendaftaran situs web\label{website_registration_design_prop}}
Sebuah situs web berupaya meningkatkan pendaftaran pengunjung pertama kali
dengan menampilkan desain situs baru kepada 1\% dari para pengunjung tersebut.
Dari 752 pengunjung yang diambil secara acak selama sebulan dan melihat
desain baru itu, 64 orang mendaftar.
\begin{parts}
\item
    Periksa setiap syarat yang diperlukan untuk membuat interval
    kepercayaan.
\item
    Hitunglah galat bakunya.
\item
    Buat dan tafsirkan interval kepercayaan 90\% bagi
    proporsi pengunjung pertama kali situs tersebut yang akan mendaftar
    dengan desain baru
    (dengan mengasumsikan perilaku pengunjung baru stabil dari waktu ke waktu).
\end{parts}
}{}

% 14

\eoce{\qt{Kupon yang mendorong kunjungan\label{store_coupon_prop}}
Sebuah toko mengambil sampel acak sebanyak 603 pembeli selama satu tahun
dan mendapati bahwa 142 di antaranya berkunjung karena kupon
yang mereka terima melalui pos.
Buatlah interval kepercayaan 95\% bagi proporsi seluruh pembeli
selama tahun tersebut yang berkunjung karena kupon yang mereka terima
melalui pos.
}{}
